% !TeX root = ../../Book3.tex
% 纲要标题：## 第2章　模的局部化、局部环与局部判据
% 纲要位置：工作记录/计划纲要.md，第 1941 行。

\chapter{模的局部化、局部环与局部判据}
\label{b3:ch:02}
\BookChapterNumber{b3:ch:02}

\begin{chapterabstract}
把允许的分母作用在模上，既会消去部分元素，也能把全局问题分解到各个素理想处。本章构造模的局部化，证明它保持正合序列，并据此逐点检测元素的消失、子模成员关系以及给定同态的单射性与满射性。有限生成与有限展示分别控制生成元和关系的共同分母；相容的局部元素则能在有限基本开覆盖上黏合。
\end{chapterabstract}

\begin{learninggoals}
\begin{itemize}
\item 用分式零判据处理模局部化中的核、商与子模。
\item 区分乘法集饱和、理想幂饱和、局部模及剩余向量空间。
\item 正确使用支集与 Hom 的局部化公式，并指出所需的有限性条件。
\item 用局部环检测给定同态，说明从点处结论扩展到邻域需要哪些有限性条件。
\item 对有限基本开覆盖检验局部模元素的相容性，并由此构造唯一的全局元素。
\end{itemize}
\end{learninggoals}

取底环 \(R=\mathbb Z\)，把 \(M=\mathbb Z/12\mathbb Z\) 看成 \(R\) 模。元素 \(\bar3\) 被 \(4\) 消去，因为 \(4\bar3=0\)。若在 \(R\) 的素理想 \((5)\) 处局部化，\(4\) 可以作为分母，\(\bar3\) 于是成为零；在 \((2)\) 处，允许的分母都是奇数，而零化 \(\bar3\) 的整数恰是 \(4\) 的倍数，所以 \(\bar3/1\) 仍非零。这里的素理想取自底环 \(R\)，不是从商环 \(\mathbb Z/12\mathbb Z\) 中选取。元素是否消失，由是否存在一个允许的分母将它消去来判定。

\ProseParagraphBreak

从一个元素转向整个模，量词便发生变化：每个生成元分别有消去它的分母，未必有一个分母同时消去所有元素。若模有限生成，把有限多个分母相乘就能得到这样的共同分母；控制全部关系则常需有限展示。支集记录哪些点仍看见这个模。另一方面，若一个给定同态在各点处都单射或都满射，局部化的正合性就能把这些结论还原为全局结论。

\ProseParagraphBreak

本章需要第一章的环局部化和素谱，以及第二卷第2章的模、核与商、第6章的张量积。除明确说明外，不假设底环是 Noether 环，也不假设所讨论的模有限生成。给单个元素清分母，与给全部生成元或全部关系选取共同分母，是两种不同的任务；这项区别决定了多数定理的有限性假设。

\ProseParagraphBreak
进一步阅读可参阅 Matsumura 的《Commutative Ring Theory》\cite[Theorems 4.4--4.6, pp. 26--27]{Matsumura1989CommutativeRingTheory}，对照局部化、张量积与局部消失的证明。该书\cite[Theorem 4.10, pp. 28--29]{Matsumura1989CommutativeRingTheory}还证明：对有限生成模及固定整数 \(r\ge0\)，局部模能由至多 \(r\) 个元素生成的点组成开集；对有限展示模，局部模为自由模的点也组成开集。

% !TeX root = ../../Book3.tex
% 纲要标题：### 2.1 模局部化
% 纲要位置：工作记录/计划纲要.md，第 1844 行。
% 纲要细目（写作范围）：
% * 分式模的构造
% * 与环局部化的张量描述
% * 局部化的正合性

\section{模局部化}
\label{b3:sec:0201}

环的局部化允许把指定元素当作分母；模也可以作同样的变换。不过，模中通常没有乘法，分式相等必须用标量作用来定义。以下始终设 \(R\) 为交换含幺环，\(S\subseteq R\) 含 \(1\) 且对乘法封闭，\(M\) 为幺性 \(R\) 模。允许 \(0\in S\)，此时局部化环及局部化模都为零。第二卷命题6.4.5已给出模局部化的张量描述；本节重新展开分母的处理，使后续局部判据能够直接使用。

\subsection{分式模的构造}
\label{b3:subsec:0201:01}

\begin{definition}\label{b3:def:module-localization}
设 \(R\) 为交换含幺环，\(S\subseteq R\) 为含一的乘法集，\(M\) 为 \(R\) 模。在 \(M\times S\) 上规定
\[
(m,s)\sim(n,t)
\quad\Longleftrightarrow\quad
\text{存在 }u\in S\text{ 使 }u(tm-sn)=0.
\]
所得等价类记为 \(m/s\)，全部等价类组成的集合记为 \(S^{-1}M\)，称为 \(M\) 关于 \(S\) 的\term[mo-jubuhua]{局部化}[localization of a module]。若 \(\mathfrak p\) 是 \(R\) 的素理想，则 \(S=R\setminus\mathfrak p\) 时记为 \(M_{\mathfrak p}\)；若 \(f\in R\)，则 \(S=\{1,f,f^2,\ldots\}\) 时记为 \(M_f\)。
\end{definition}
\symidx{\(S^{-1}M\)、\(M_{\mathfrak p}\)、\(M_f\)：模的局部化}

先核对这确实是等价关系。自反性取 \(u=1\)，对称性把等式取负。若 \(u(tm-sn)=0\)、\(v(rn-tp)=0\)，则
\[
uvt(rm-sp)=uv\bigl(r(tm-sn)+s(rn-tp)\bigr)=0,
\]
所以传递性也成立。关系中的额外乘子 \(u\) 不能省去：若某个分母已经零化了 \(m\)，局部化后 \(m\) 就应消失。

\begin{proposition}\label{b3:prop:module-fraction-operations}
设 \(R\) 为交换含幺环，\(S\subseteq R\) 为含一的乘法集，\(M\) 为 \(R\) 模。下列公式使 \(S^{-1}M\) 成为 \(S^{-1}R\) 模：
\[
\frac ms+\frac nt=\frac{tm+sn}{st},\qquad
\frac as\frac mt=\frac{am}{st}.
\]
自然映射 \(\iota_M:M\rightarrow S^{-1}M\)、\(m\mapsto m/1\) 为 \(R\) 线性映射，且
\begin{equation}\label{b3:eq:localized-element-zero}
\frac ms=0\quad\Longleftrightarrow\quad
um=0\text{ 对某个 }u\in S\text{ 成立}.
\end{equation}
若 \(N\) 是 \(S^{-1}R\) 模，任意 \(R\) 线性映射 \(h:M\rightarrow N\) 唯一延拓为 \(S^{-1}R\) 线性映射
\[
\widetilde h:S^{-1}M\longrightarrow N,\qquad
\widetilde h(m/s)=(s/1)^{-1}h(m).
\]
这唯一的 \(S^{-1}R\) 线性延拓使下图交换：
\[
\begin{tikzcd}[column sep=large,row sep=large]
M\arrow[r,"\iota_M"]\arrow[dr,"h"']
&S^{-1}M\arrow[d,dashed,"\widetilde h"]\\
&N
\end{tikzcd}
\]
\end{proposition}
\begin{proof}
若 \(m/s=m'/s'\)，取 \(u\in S\) 使 \(u(s'm-sm')=0\)。加上同一个 \(n/t\) 后，两分子的交叉差为 \(t^2(s'm-sm')\)，仍被 \(u\) 零化；乘上同一个 \(a/t\) 后，交叉差为 \(at(s'm-sm')\)，也被 \(u\) 零化。改变第二个加数的代表元同理。若 \(a/t=a'/t'\)，相应交叉差为 \(s(t'a-ta')m\)，由环分式的等价关系零化。因此运算良定义。把有限个分式通分后，结合律、分配律、零元 \(0/1\)、负元 \(-m/s\) 以及单位作用，都化为 \(M\) 中的相应公理。零判据直接把 \(m/s\) 与 \(0/1\) 比较。

若 \(u(tm-sn)=0\)，在 \(N\) 中 \(u,s,t\) 均可逆，施加 \(h\) 再消去这些单位，即得 \(s^{-1}\cdot h(m)=t^{-1}\cdot h(n)\)。这证明延拓良定义；通分验证线性。因 \(m/s=(1/s)(m/1)\)，任何这样的延拓都只能采用所给公式，故唯一。
\end{proof}

例如 \(R=\mathbb Z\)、\(M=\mathbb Z/12\mathbb Z\)，在 \((2)\) 处局部化时 \(3\) 可逆，原来阶为 \(3\) 的部分消失，得到 \(M_{(2)}\cong\mathbb Z/4\mathbb Z\)。确切的同构由 \(\bar a/s\mapsto \bar a\bar s^{-1}\) 给出，其中 \(s\) 为奇数；其核为零可由 \(4\mid a\Rightarrow12\mid3a\) 核验。

\subsection{模局部化的张量描述}
\label{b3:subsec:0201:02}

\begin{proposition}\label{b3:prop:localization-tensor}
设 \(R\) 为交换含幺环，\(S\subseteq R\) 为含一的乘法集，\(M\) 为 \(R\) 模。存在关于 \(M\) 自然的 \(S^{-1}R\) 模同构
\[
(S^{-1}R)\otimes_R M\longrightarrow S^{-1}M,
\qquad \frac as\otimes m\longmapsto\frac{am}s.
\]
其逆映射为 \(m/s\mapsto(1/s)\otimes m\)。
\end{proposition}
\begin{proof}
正向公式对两个变量可加，且 \((ar/s)\otimes m\) 与 \((a/s)\otimes rm\) 具有相同的像，故由张量积的泛性质给出线性映射。若 \(u(tm-sn)=0\)，则
\[
\frac1s\otimes m-\frac1t\otimes n
=\frac1{ust}\otimes u(tm-sn)=0.
\]
因此逆向公式良定义，通分后可直接验证它线性。两种公式的复合在 \(m/s\) 及纯张量 \((a/s)\otimes m\) 上均为恒等，故互为逆映射。模同态 \(f:M\rightarrow N\) 在两边都把分子 \(m\) 变为 \(f(m)\)，这就是自然性。
\end{proof}

局部化因而是沿 \(R\rightarrow S^{-1}R\) 的标量扩张。它允许把原来的生成元乘以新的分式标量，并消去被某个分母零化的元素。例如对任意理想 \(I\)，该描述给出
\[
S^{-1}(M/IM)\cong(S^{-1}M)/(S^{-1}I)(S^{-1}M).
\]
这里也可以把分式 \((m+IM)/s\) 直接送到 \(m/s\) 的商类；其良定义和核将在下一小节的正合性中一并得到。

\subsection{局部化的正合性}
\label{b3:subsec:0201:03}

\begin{theorem}\label{b3:thm:localization-exact}
设 \(R\) 为交换含幺环，\(S\subseteq R\) 为含一的乘法集，\(L,M,N\) 为 \(R\) 模。若 \(f:M\rightarrow N\) 是 \(R\) 模同态，则它诱导映射 \(S^{-1}f(m/s)=f(m)/s\)。这一定义给出正合函子：若
\[
0\longrightarrow L\xrightarrow{i}M\xrightarrow{q}N\longrightarrow0
\]
正合，则局部化后的序列仍正合。因此局部化与核、像、余核及有限交相容。
\end{theorem}
\begin{proof}
若两分式相等，对它们的交叉差施加 \(f\) 即可证明诱导映射良定义；恒等映射与复合的相容性在分子上直接成立。

若 \(i(l)/s=0\)，由\eqref{b3:eq:localized-element-zero}有 \(u\cdot i(l)=i(ul)=0\)。原映射 \(i\) 单射，故 \(ul=0\)，于是 \(l/s=0\)。若 \(q(m)/s=0\)，取 \(u\in S\) 使 \(q(um)=0\)，再取 \(l\in L\) 使 \(i(l)=um\)。则
\[
\frac ms=\frac{um}{us}=S^{-1}i\left(\frac l{us}\right).
\]
反过来，\(qi=0\) 保证像包含于核。最后，对 \(n/s\) 选 \(m\) 满足 \(q(m)=n\)，便有原像 \(m/s\)，证明满射性。

把任意同态分别分解为到其像的满射及像的包含，便得到关于核、像和余核的断言。对子模 \(L_1,L_2\subseteq M\)，把交写为映射 \(L_1\oplus L_2\rightarrow M\)、\((x,y)\mapsto x-y\) 的核，再用 \(S^{-1}(L_1\oplus L_2)\cong S^{-1}L_1\oplus S^{-1}L_2\)，得到有限交的相容性。
\end{proof}

第二卷命题7.3.10用有限关系判据证明 \(S^{-1}R\) 是平坦 \(R\) 模；结合\cref{b3:prop:localization-tensor}，那也能推出正合性。这里的分式证明则明确指出：落在局部核中不意味着原分子已经在核中，而是要先乘上某个分母。这个区别将在理想的收缩和模的支集中反复使用。

% !TeX root = ../../Book3.tex
% 纲要标题：### 2.2 理想与局部化
% 纲要位置：工作记录/计划纲要.md，第 1949 行。
% 纲要细目（写作范围）：
% * 扩张、收缩和饱和；有向并与单元素饱和的几何实例
% * 局部化后的素理想对应
% * 剩余域与纤维；局部模和剩余向量空间、空纤维及非约化纤维

\section{理想与局部化}
\label{b3:sec:0202}

局部化可能使一个非零理想变成零，也可能使一个真理想变成单位理想。要准确追踪理想，必须同时使用扩张与收缩；回到原环后增加的元素，正是能被某个允许的分母送入原理想的元素。本节还把第一章的素谱对应用于剩余域及纤维。

\subsection{扩张、收缩与饱和}
\label{b3:subsec:0202:01}

设 \(R\) 为交换含幺环，\(S\subseteq R\) 为含一的乘法集，\(I\) 为 \(R\) 的理想。由\cref{b3:thm:localization-exact}，包含映射 \(I\hookrightarrow R\) 局部化后仍为单射，因而可把模的局部化 \(S^{-1}I\) 识别为 \(S^{-1}R\) 中由分式 \(a/s\)、\(a\in I\)、\(s\in S\) 组成的子模。这些分式也对乘以任意环分式封闭，而且每个 \(a/s=(a/1)(1/s)\)，所以其像正是扩张理想 \(I(S^{-1}R)\)。以下用 \(S^{-1}I\) 同时表示这个扩张理想。对 \(S^{-1}R\) 的理想 \(J\)，其收缩是 \(J^c=\{r\in R:r/1\in J\}\)。

\begin{definition}\label{b3:def:multiplicative-saturation}
设 \(R\) 为交换含幺环，\(S\subseteq R\) 为含一的乘法集，\(I\) 为 \(R\) 的理想。对 \(s\in R\)，记 \((I:s)=\{\,r\in R\mid sr\in I\,\}\)。\(I\) 关于 \(S\) 的\term[chengfaji-baohe]{饱和}[saturation with respect to a multiplicative set]为
\[
\operatorname{Sat}_S(I)\coloneqq
\{\,r\in R\mid sr\in I\text{ 对某个 }s\in S\text{ 成立}\,\}
=\bigcup_{s\in S}(I:s).
\]
若 \(\operatorname{Sat}_S(I)=I\)，则称 \(I\) 关于 \(S\) 饱和。
\end{definition}
\symidx{\(\operatorname{Sat}_S(I)\)：理想关于乘法集的饱和}

理想的任意并一般不是理想，但这里的并具有额外的相容性。对 \(s,t\in S\)，有
\[
(I:s)\subseteq(I:st),\qquad (I:t)\subseteq(I:st).
\]
因为 \(st\in S\)，任意两个成员都包含于这族理想中的同一个成员，即这族理想在包含关系下有向。因此并中的任意两个元素可以在同一理想内相加，每个元素乘以 \(R\) 的任意标量也仍在并中，故 \(\operatorname{Sat}_S(I)\) 确实是理想。

\begin{proposition}\label{b3:prop:ideal-localization-saturation}
设 \(R\) 为交换含幺环，\(S\subseteq R\) 为含一的乘法集，\(I\) 为 \(R\) 的理想，\(J\) 为 \(S^{-1}R\) 的理想。记 \(S^{-1}I\) 为扩张理想，\(J^c=\{r\in R:r/1\in J\}\) 为收缩理想。则
\[
(S^{-1}I)^c=\operatorname{Sat}_S(I),\qquad S^{-1}(J^c)=J.
\]
因此 \(S^{-1}R\) 的理想与 \(R\) 中关于 \(S\) 饱和的理想一一对应。此外，\(S^{-1}I=S^{-1}R\) 当且仅当 \(I\cap S\ne\varnothing\)。
\end{proposition}
\begin{proof}
若 \(r/1=a/s\)、\(a\in I\)，则存在 \(u\in S\) 使 \(u(sr-a)=0\)，所以 \((us)r=ua\in I\)。反过来，若 \(sr\in I\)，则 \(r/1=(sr)/s\in S^{-1}I\)。这证明第一个等式。

对 \(a/s\in J\)，有 \(a/1=(s/1)(a/s)\in J\)，因而 \(a\in J^c\)；反向包含由 \(J\) 是理想得到。两个等式还表明扩张收缩在饱和理想上互逆。最后，\(1\in(S^{-1}I)^c\) 等价于存在 \(s\in S\) 属于 \(I\)，这就是单位理想的判据。
\end{proof}

设 \(k\) 为域。在 \(R=k[x,y]\) 中，取 \(S=\{1,x,x^2,\ldots\}\)、\(I=(xy)\)，则
\[
\operatorname{Sat}_S(I)=(y),\qquad I R_x=(y)R_x.
\]
确实，\(x(y)\subseteq I\) 给出一个包含；若 \(x^n f\in(xy)\)，在整环 \(k[x,y]/(y)=k[x]\) 中约去非零的 \(x^n\)，便得 \(f\in(y)\)。同一个 \(I\) 在 \(x,y\) 都可逆的环中扩张成单位理想。

\ProseParagraphBreak
在这个例子中，\(V(xy)=V(x)\cup V(y)\)。把 \(x\) 变成单位相当于限制到 \(D(x)\)，其中不再有 \(V(x)\) 上的点，剩下的是 \(V(y)\cap D(x)\)。这个集合包含 \(V(y)\) 的一般点 \((y)\)，因为 \(x\notin(y)\)；所以它在原谱中的闭包正是 \(V(y)\)。所算出的饱和理想 \((y)\) 是根理想，恰好对应这个闭包。

\ProseParagraphBreak
若 \(S\) 由单个 \(f\in R\) 的幂组成，则 \(\operatorname{Sat}_S(I)=I:f^\infty=\bigcup_{n\ge0}(I:f^n)\)。它与关于多生成元理想 \(J\) 的 \(I:J^\infty\) 不同：后者要求 \(J^n r\subseteq I\)，即该幂理想的每个元素都满足条件；乘法集的饱和只要求存在一个分母。理想商与饱和的系统计算将在\chapref{b3:ch:08}继续讨论。

\subsection{局部化中的素理想对应}
\label{b3:subsec:0202:02}

\begin{proposition}\label{b3:prop:localized-prime-ideals}
设 \(R\) 为交换含幺环，\(S\subseteq R\) 为含一的乘法集。扩张与收缩给出素理想之间的对应
\[
\{\,\mathfrak p\in\operatorname{Spec}R\mid\mathfrak p\cap S=\varnothing\,\}
\longleftrightarrow\operatorname{Spec}(S^{-1}R),
\qquad \mathfrak p\longmapsto S^{-1}\mathfrak p.
\]
对 \(R\) 的任一素理想 \(\mathfrak p\)，在 \(R_{\mathfrak p}\) 中唯一的极大理想为 \(\mathfrak pR_{\mathfrak p}\)，其素理想对应于 \(R\) 中包含在 \(\mathfrak p\) 内的素理想。
\end{proposition}
\begin{proof}
第一章的\cref{b3:thm:spec-localization}已证明这一对应及其拓扑相容性；这里说明代数上怎样使用。若 \(sr\in\mathfrak p\)、\(s\notin\mathfrak p\)，素性给出 \(r\in\mathfrak p\)，故 \(\mathfrak p\) 饱和。商环
\[
(S^{-1}R)/(S^{-1}\mathfrak p)\cong\bar S^{-1}(R/\mathfrak p)
\]
是非零整环，因此扩张仍为素理想。相反，局部化环中所有 \(s/1\) 都是单位，不可能落入真素理想；其收缩是与 \(S\) 不交的素理想。由\cref{b3:prop:ideal-localization-saturation}两种操作互逆。

取 \(S=R\setminus\mathfrak p\)。一个素理想与 \(S\) 不交正是它包含于 \(\mathfrak p\)，故最大的一个为 \(\mathfrak p\)。也可直接看出，若 \(a/s\notin\mathfrak pR_{\mathfrak p}\)，则 \(a\notin\mathfrak p\)，其逆为 \(s/a\)；所以剩下的全部元素都是单位，极大理想唯一。
\end{proof}

\subsection{剩余域与纤维}
\label{b3:subsec:0202:03}

设 \(R,A\) 为交换含幺环。保单位环同态 \(\varphi:R\rightarrow A\) 把 \(A\) 的素理想 \(\mathfrak q\) 送到收缩 \(\varphi^{-1}(\mathfrak q)\)。固定 \(R\) 的一个素理想 \(\mathfrak p\)，哪些点恰好落在它上面？要使收缩等于 \(\mathfrak p\)，可以先把 \(\varphi(\mathfrak p)\) 中的元素模掉，再把 \(\varphi(R\setminus\mathfrak p)\) 中的元素变成单位。这两种操作合起来，就是从 \(R\) 到 \(\mathfrak p\) 处的剩余域作标量变换。

\begin{definition}\label{b3:def:residue-field-fibre}
设 \(R\) 为交换含幺环，\(\mathfrak p\) 为 \(R\) 的素理想。\(\mathfrak p\) 处的剩余域为
\[
\kappa(\mathfrak p)\coloneqq R_{\mathfrak p}/\mathfrak pR_{\mathfrak p}
\cong\operatorname{Frac}(R/\mathfrak p).
\]
给定交换含幺环 \(A\) 及保单位环同态 \(\varphi:R\rightarrow A\)，称 \(A\otimes_R\kappa(\mathfrak p)\) 为 \(\varphi\) 在 \(\mathfrak p\) 处的\term[xianwei-daishu]{纤维代数}[fibre algebra]。对任一 \(R\) 模 \(M\)，它在 \(\mathfrak p\) 处的剩余向量空间为 \(M\otimes_R\kappa(\mathfrak p)\cong M_{\mathfrak p}/\mathfrak pM_{\mathfrak p}\)。
\end{definition}
\symidx{\(\kappa(\mathfrak p)\)：素理想处的剩余域}

剩余域的同构可从分式直接看出：模掉 \(\mathfrak p\) 后，\(R\setminus\mathfrak p\) 的像恰是整环 \(R/\mathfrak p\) 的全部非零元素。两个张量描述则依次使用\cref{b3:prop:localization-tensor}及商模的张量公式。

\ProseParagraphBreak
局部模 \(M_{\mathfrak p}\) 只把 \(R\setminus\mathfrak p\) 中的标量变成单位；剩余向量空间还进一步模掉 \(\mathfrak pM_{\mathfrak p}\)，因此两者保留的信息不同。例如设 \(k\) 为域，取 \(R=k[t]_{(t)}\)、\(\mathfrak p=(t)R\)、\(M=R/(t^2)\)。\(R\setminus\mathfrak p\) 中的元素已经是单位，故
\[
M_{\mathfrak p}\cong R/(t^2),\qquad
M\otimes_R\kappa(\mathfrak p)\cong M\otimes_R k\cong k.
\]
在局部模中，\(t\) 的商类 \(\bar t\) 非零，却满足 \(t\bar t=0\)；在剩余向量空间中，\(\bar t\otimes1=\bar1\otimes(t\cdot1)=0\)。所以取纤维会丢掉这个局部模中仍然可见的非零元素。

\begin{proposition}\label{b3:prop:fibre-spectrum}
设 \(R,A\) 为交换含幺环，\(\varphi:R\to A\) 为保单位环同态，\(\mathfrak p\) 为 \(R\) 的素理想。纤维代数 \(A\otimes_R\kappa(\mathfrak p)\) 的素理想自然对应于 \(A\) 中满足 \(\varphi^{-1}(\mathfrak q)=\mathfrak p\) 的素理想 \(\mathfrak q\)。
\end{proposition}
\begin{proof}
令 \(S=R\setminus\mathfrak p\)。通过 \(\varphi\) 赋予 \(A\) 一个 \(R\) 代数结构，并记
\[
S^{-1}A\coloneqq\varphi(S)^{-1}A.
\]
这里 \(S^{-1}A\) 是把 \(\varphi(S)\) 中的元素变成单位所得的环；\(\mathfrak pS^{-1}A\) 表示 \(\varphi(\mathfrak p)\) 在该环中生成的理想。局部化及取商的张量公式给出
\[
A\otimes_R\kappa(\mathfrak p)
\cong S^{-1}A/\mathfrak pS^{-1}A.
\]
商环的素理想对应于 \(S^{-1}A\) 中包含 \(\mathfrak pS^{-1}A\) 的素理想；局部化对应又把它们识别为 \(A\) 中包含 \(\varphi(\mathfrak p)\) 且与 \(\varphi(S)\) 不交的素理想。这两个条件给出两边的包含：
\[
\begin{aligned}
\varphi(\mathfrak p)\subseteq\mathfrak q
&\ \Longrightarrow\ \mathfrak p\subseteq\varphi^{-1}(\mathfrak q),\\
\varphi(R\setminus\mathfrak p)\cap\mathfrak q=\varnothing
&\ \Longrightarrow\ \varphi^{-1}(\mathfrak q)\subseteq\mathfrak p.
\end{aligned}
\]
第二个方向是因为：若 \(r\in\varphi^{-1}(\mathfrak q)\) 却不属于 \(\mathfrak p\)，则 \(r\in S\)，于是 \(\varphi(r)\) 同时属于 \(\varphi(S)\) 和 \(\mathfrak q\)，矛盾。反过来，若 \(\varphi^{-1}(\mathfrak q)=\mathfrak p\)，则上述两个条件都成立。因此它们合起来恰好是所述等式。
\end{proof}

\begin{figure}[htbp]
\centering
\begin{tikzcd}[column sep=large,row sep=large]
\operatorname{Spec}\bigl(A\otimes_R\kappa(\mathfrak p)\bigr)
\arrow[r,"\iota"] \arrow[d]
& \operatorname{Spec}A \arrow[d,"\varphi^*"] \\
\{\mathfrak p\} \arrow[r,hook]
& \operatorname{Spec}R
\end{tikzcd}
\caption{谱映射在一个点上的纤维}
\label{b3:fig:spectrum-fibre}
\end{figure}

\cref{b3:fig:spectrum-fibre}中的 \(\iota\) 由 \(a\mapsto a\otimes1\) 诱导。上面的命题说明，它的像恰为 \(\{\,\mathfrak q\in\operatorname{Spec}A\mid\varphi^{-1}(\mathfrak q)=\mathfrak p\,\}\)，也就是 \(\varphi^*\) 在 \(\mathfrak p\) 上的集合论纤维；这解释了纤维代数的名称。

\ProseParagraphBreak
纤维也可能为空。例如 \(\mathbb Z\hookrightarrow\mathbb Q\) 在 \((2)\) 处的纤维代数为 \(\mathbb Q\otimes_{\mathbb Z}\mathbb F_2\)。它的单位元满足 \(1\otimes\bar1=(1/2)\otimes2\bar1=0\)，所以纤维代数是零环，素谱为空。

\ProseParagraphBreak
例如，设 \(k\) 为域，\(\varphi:k[t]\rightarrow k[x]\) 为满足 \(\varphi(t)=x^2\) 的 \(k\) 代数同态。对 \(a\in k\)，\(\varphi\) 在 \((t-a)\) 处的纤维代数为 \(k[x]/(x^2-a)\)。当 \(a=0\) 时，其中 \(x\) 的像非零却平方为零；素谱只有一个点，其剩余域为 \(k\)，但纤维代数仍保存这项幂零信息。当 \(a=b^2\ne0\)、\(b\in k\) 且 \(\operatorname{char}k\ne2\) 时，\(x-b\) 与 \(x+b\) 互素，中国剩余定理给出 \(k[x]/(x^2-a)\cong k\times k\)，有两个点，剩余域都是 \(k\)。若 \(\operatorname{char}k=2\) 且 \(a=b^2\ne0\)，则 \(x^2-a=(x-b)^2\)；令 \(\varepsilon\) 为 \(x-b\) 的商类，纤维代数同构于 \(k[\varepsilon]/(\varepsilon^2)\)。这仍只有一个剩余域为 \(k\) 的点，且 \(\varepsilon\ne0\) 是幂零元，即一个非约化双重点。若 \(a\) 在 \(k\) 中没有平方根，则 \(x^2-a\) 不可约，纤维代数是二次域扩张，只有一个点，其剩余域就是该二次扩张，而且没有非零幂零元。这些情形说明，点数、点的剩余域和纤维代数的约化性是三种不同的信息；同样只有一个点，既可能有幂零元，也可能是域。

% !TeX root = ../../Book3.tex
% 纲要标题：### 2.3 模的支集
% 纲要位置：工作记录/计划纲要.md，第 1955 行。
% 纲要细目（写作范围）：
% * 支集、零化子与有限生成性；循环模支集等于理想闭集
% * 局部化与 Hom 的交换条件；有限生成不能代替有限展示的循环模反例
% * 支集与短正合列

\section{模的支集}
\label{b3:sec:0203}

局部化把一个模分散到素谱的各点；支集记录它在哪些点仍然非零。有限生成性使这些局部信息受到一个理想的统一控制。处理同态时，还要进一步区分“有限个生成元”与“有限个生成关系”，因为清分母必须一次覆盖全部有关数据。

\subsection{支集、零化子与有限生成性}
\label{b3:subsec:0203:01}

\begin{definition}\label{b3:def:support-annihilator}
设 \(R\) 为交换含幺环，\(M\) 为 \(R\) 模。\(M\) 的\term[zhiji]{支集}[support]与\term[linghuazi]{零化子}[annihilator]分别为
\[
\begin{aligned}
\operatorname{Supp}_R M
&=\{\,\mathfrak p\in\operatorname{Spec}R\mid M_{\mathfrak p}\ne0\,\},\\
\operatorname{Ann}_R M
&=\{\,a\in R\mid am=0\text{ 对所有 }m\in M\,\}.
\end{aligned}
\]
零化子也称湮灭子。\termidx[yinmiezi]{湮灭子}对单个元素，简记
\(\operatorname{Ann}_R(m)=\{a\in R:am=0\}\)；它等于循环子模 \(Rm\) 的零化子。
\end{definition}
\symidx{\(\operatorname{Supp}_R M\)：模的支集}
\symidx{\(\operatorname{Ann}_R M\)：模的零化子}

\begin{theorem}\label{b3:thm:support-finite}
设 \(R\) 为交换含幺环，\(M\) 为 \(R\) 模。则
\[
\operatorname{Supp}_R M\subseteq V(\operatorname{Ann}_R M).
\]
若 \(M\) 有限生成，则等号成立。更具体地，对每个 \(R\) 的素理想 \(\mathfrak p\)，有限生成模 \(M\) 满足
\[
M_{\mathfrak p}=0
\quad\Longleftrightarrow\quad
sM=0\text{ 对某个 }s\notin\mathfrak p\text{ 成立}.
\]
\end{theorem}
\begin{proof}
若 \(a\in\operatorname{Ann}_R M\setminus\mathfrak p\)，在 \(R_{\mathfrak p}\) 中 \(a\) 可逆，又零化 \(M_{\mathfrak p}\)，故该模为零。这证明一般包含关系。

设 \(M=Rm_1+\cdots+Rm_r\) 且 \(M_{\mathfrak p}=0\)。由\eqref{b3:eq:localized-element-zero}，对每个 \(i\) 有 \(s_i\notin\mathfrak p\) 使 \(s_im_i=0\)。取 \(s=s_1\cdots s_r\)，则 \(s\notin\mathfrak p\) 且 \(sM=0\)。反向由 \(s\) 局部可逆得到。于是 \(M_{\mathfrak p}=0\) 等价于 \(\operatorname{Ann}_R M\not\subseteq\mathfrak p\)，即所述等式。对零模取空生成集及 \(s=1\)，论证仍成立。
\end{proof}

因而有限生成模的支集是闭集。特别地，对 \(R\) 的任意理想 \(I\)，循环模 \(R/I\) 有限生成，且 \(\operatorname{Ann}_R(R/I)=I\)：零化它的元素必须零化 \(1+I\)，反向则由商模的运算得到。因此
\[
\operatorname{Supp}_R(R/I)=V(I).
\]
例如，对域 \(k\)，有
\[
\operatorname{Supp}_{k[x,y]} k[x,y]/(xy)=V(xy)=V(x)\cup V(y).
\]
模掉 \(xy\) 得到的模同时分布在两条坐标轴上；若进一步在 \(x\) 处局部化，则只留下 \(y=0\) 且 \(x\ne0\) 的部分。

\ProseParagraphBreak
有限生成条件不能省去。整数模 \(\mathbb Q/\mathbb Z\) 的零化子为零：对非零整数 \(n\)，选不能整除 \(n\) 的正整数 \(d\)，则 \(n(1/d+\mathbb Z)\ne0\)。但在 \((0)\) 处局部化后，每个元素都被某个非零整数零化，故整个模变成零。另一方面，对任一素数 \(p\)，元素 \(1/p+\mathbb Z\) 在 \((p)\) 处局部化后仍不为零：若它变成零，就有某个 \(s\in\mathbb Z\setminus(p)\) 使 \(s/p\in\mathbb Z\)，与 \(p\nmid s\) 矛盾。由于 \(\mathbb Z\) 的素理想只有 \((0)\) 及各个 \((p)\)，这就完整算出
\[
\operatorname{Supp}_{\mathbb Z}(\mathbb Q/\mathbb Z)
=\{\,(p)\mid p\text{ 为素数}\,\},\qquad
V\bigl(\operatorname{Ann}_{\mathbb Z}(\mathbb Q/\mathbb Z)\bigr)
=\operatorname{Spec}\mathbb Z.
\]
非零整数不可能被所有素数整除，故 \(\bigcap_{p\text{ 为素数}}(p)=(0)\)。若闭集 \(V(I)\) 包含上述支集，则 \(I\) 包含于这个交，只能为零。因此支集在 \(\operatorname{Spec}\mathbb Z\) 中稠密；它又缺少 \((0)\)，因而不是闭集。此例中每个元素各有一个零化它的非零乘子，却没有一个非零整数同时零化全部元素。

\subsection{局部化与 Hom 的交换条件}
\label{b3:subsec:0203:02}

设 \(R\) 为交换含幺环，\(S\subseteq R\) 为含一的乘法集，\(M,N\) 为 \(R\) 模。局部化后的同态可以对分式取值，因而有自然映射
\begin{equation}\label{b3:eq:hom-localization-map}
S^{-1}\operatorname{Hom}_R(M,N)
\longrightarrow
\operatorname{Hom}_{S^{-1}R}(S^{-1}M,S^{-1}N),
\qquad
\frac fs\longmapsto
\left(\frac mt\longmapsto\frac{f(m)}{st}\right).
\end{equation}
若两侧的分式代表元改变，交叉差仍被某个分母零化；因此该公式良定义。也可先在 \(M\) 上取 \(m\mapsto f(m)/s\)，再用\cref{b3:prop:module-fraction-operations}的延拓泛性质。

\begin{theorem}\label{b3:thm:hom-localization}
设 \(R\) 为交换含幺环，\(S\subseteq R\) 为含一的乘法集，\(M,N\) 为 \(R\) 模。考虑自然映射
\[
S^{-1}\operatorname{Hom}_R(M,N)\longrightarrow
\operatorname{Hom}_{S^{-1}R}(S^{-1}M,S^{-1}N),\qquad
f/s\longmapsto\bigl(m/t\longmapsto f(m)/(st)\bigr).
\]
若 \(M\) 有限生成，则此映射为单射。若 \(M\) 有限展示，即存在非负整数 \(a,b\) 及 \(R\) 模正合列
\[
R^a\longrightarrow R^b\longrightarrow M\longrightarrow0,
\]
则该映射为同构。两种结论都不要求 \(N\) 有限生成。
\end{theorem}
\begin{proof}
先只假设 \(M\) 有限生成，取生成元 \(m_1,\ldots,m_r\)。单射性要求：一个同态若在局部化后消失，就能找到一个分母零化这个同态的全部值。若 \(f/s\) 在右侧为零，则每个 \(f(m_i)/s=0\)，故可取 \(u_i\in S\) 使 \(u_i\cdot f(m_i)=0\)。因为生成元只有有限个，乘积 \(u=\prod_{i=1}^r u_i\) 零化 \(f\) 在全部生成元上的值，因而 \(uf=0\)，即 \(f/s=0\)。这证明单射性；当 \(M=0\) 时可取空生成集及 \(u=1\)。

再设 \(M\) 有上述有限展示。要把局部化后的同态写成 \(f/s\)，仅给生成元的像通分还不够；通分后的像必须在 \(N\) 中满足全部关系，才能定义 \(f:M\to N\)。记 \(e_1,\ldots,e_b\) 在 \(M\) 中的像为 \(m_i\)。展示的第一个映射的像由有限个关系
\(\sum_i a_{ij}m_i=0\)、\(1\le j\le a\) 生成。给定右侧的同态 \(g\)，把有限个 \(g(m_i/1)\) 通分为 \(n_i/s\)。每个关系给出
\[
\frac{\sum_i a_{ij}n_i}{s}=0.
\]
故可选一个共同的 \(u\in S\)，使全部 \(\sum_i a_{ij}u n_i=0\) 在 \(N\) 中成立。于是 \(m_i\mapsto u n_i\) 满足全部关系，诱导 \(f:M\rightarrow N\)。此时 \(f/(us)\) 在\eqref{b3:eq:hom-localization-map}下的像在 \(m_i/1\) 上等于 \(g\)，因而就是 \(g\)。这证明满射性。
\end{proof}

有限生成不能代替有限展示来保证满射性。下面的循环模只有一个生成元，却无法同时清除所有关系中的分母。

\begin{proposition}\label{b3:prop:hom-localization-fg-counterexample}
设 \(k\) 为域，\(R=\prod_{n\ge1}k\)，\(I=\bigoplus_{n\ge1}k\subseteq R\) 为只有有限多个非零坐标的元素组成的理想。对有限集 \(F\subseteq\mathbb Z_{\ge1}\)，令 \(e_F\) 在 \(F\) 的坐标上等于 \(1\)，其余坐标上等于 \(0\)，并取 \(S=\{1-e_F:F\text{ 有限}\}\)。则 \(S\) 是含一而不含零的乘法集，\(M=R/I\) 是有限生成但非有限展示的 \(R\) 模，且自然映射
\[
S^{-1}\operatorname{Hom}_R(M,R)\longrightarrow
\operatorname{Hom}_{S^{-1}R}(S^{-1}M,S^{-1}R)
\]
的源为零，目标非零。还有自然同构 \(S^{-1}R\cong R/I\)，把 \(a/s\) 送到 \(a+I\)。
\end{proposition}
\begin{proof}
因为 \((1-e_F)(1-e_G)=1-e_{F\cup G}\)，\(S\) 对乘法封闭；取空集得到 \(1\in S\)，而有限集不可能包含所有正整数，所以 \(0\notin S\)。\(M\) 由 \(1+I\) 生成。一个同态 \(M\to R\) 在该生成元上的值 \(a\) 必须满足 \(Ia=0\)。用每个单坐标幂等元 \(e_{\{n\}}\in I\) 相乘得到 \(a_n=0\)，故 \(a=0\)，即
\[
\operatorname{Hom}_R(M,R)=\operatorname{Ann}_R I=0.
\]
另一方面，每个 \(x\in I\) 只有有限多个非零坐标；若这些坐标包含于 \(F\)，则 \((1-e_F)x=0\)。由分式零判据，\(S^{-1}I=0\)，局部化正合性于是给出 \(S^{-1}M\cong S^{-1}R\)。环 \(S^{-1}R\) 非零，因为 \(1/1=0\) 将要求 \(S\) 含有零元。因此目标包含非零恒等同态，而源为零，映射不满射。由\cref{b3:thm:hom-localization}，\(M\) 不可能有限展示。

对每个 \(s=1-e_F\in S\)，其在 \(R/I\) 中的像等于 \(1\)，故商映射延拓为 \(S^{-1}R\to R/I\)、\(a/s\mapsto a+I\)。其逆为 \(a+I\mapsto a/1\)：\(S^{-1}I=0\) 保证该映射良定义，且 \(e_F/1=0\) 使 \(s/1=1\)，从而 \(a/s=a/1\)。
\end{proof}

若连有限生成也不假设，取 \(R=\mathbb Z\)、\(S=\mathbb Z\setminus\{0\}\)、\(M=\mathbb Q\)、\(N=\mathbb Z\)。对任意 \(\mathbb Z\) 模同态 \(f:\mathbb Q\rightarrow\mathbb Z\)，有
\[
f(q)=n\cdot f(q/n)\qquad(q\in\mathbb Q,\ n\ge1).
\]
所以整数 \(f(q)\) 被每个正整数整除，只能为零。因而 \(\operatorname{Hom}_{\mathbb Z}(\mathbb Q,\mathbb Z)=0\)，而局部化后两模都成为 \(\mathbb Q\)，右侧为 \(\operatorname{Hom}_{\mathbb Q}(\mathbb Q,\mathbb Q)=\mathbb Q\)。因此一般的 Hom 不能无条件移过局部化。

\subsection{短正合列中的支集}
\label{b3:subsec:0203:03}

\begin{proposition}\label{b3:prop:support-exact}
设 \(R\) 为交换含幺环，\(L,M,N\) 为 \(R\) 模。若 \(0\rightarrow L\rightarrow M\rightarrow N\rightarrow0\) 正合，则
\[
\operatorname{Supp}_R M
=\operatorname{Supp}_R L\cup\operatorname{Supp}_R N.
\]
该结论对任意 \(R\) 模成立。此外，任意一族 \(R\) 模 \((M_\lambda)_\lambda\) 满足
\[
\operatorname{Supp}_R\left(\bigoplus_{\lambda}M_\lambda\right)
=\bigcup_{\lambda}\operatorname{Supp}_R M_\lambda.
\]
\end{proposition}
\begin{proof}
由\cref{b3:thm:localization-exact}，在每个 \(\mathfrak p\) 处仍有短正合列。中间项为零当且仅当两端同时为零，取补集即得第一式。

局部化与直和相容：把分式 \((m_\lambda)_\lambda/s\) 送到 \((m_\lambda/s)_\lambda\)。其像仍只有有限多个非零分量；反过来，对这些非零分量的分母取乘积，就能通分并得到逆映射。这里的有限性来自直和的定义，与素谱上模的支集是不同的概念。直和为零当且仅当各分量为零，因而得到第二式。
\end{proof}

第一式说明支集忽略了短正合列是否分裂。例如，设 \(k\) 为域，在 \(R=k[x]\) 上有
\[
0\longrightarrow(x)/(x^2)\longrightarrow R/(x^2)
\longrightarrow R/(x)\longrightarrow0
\]
的三个非零模支集都为 \(\bigl\{(x)\bigr\}\)，但中间项中的 \(x\) 作用非零而平方为零，两端的 \(x\) 作用均为零。这种厚度信息并未由支集单独记录，后面还要用长度、关联素理想及准素分解作更细的分析。

% !TeX root = ../../Book3.tex
% 纲要标题：### 2.4 局部—整体原则
% 纲要位置：工作记录/计划纲要.md，第 1862 行。
% 纲要细目（写作范围）：
% * 元素和模的消失判据
% * 单射、满射及同构的局部检测
% * 有限基本开覆盖与邻域有限性
% * 有限基本开覆盖上的相容局部元素黏合

\section{局部—整体原则}
\label{b3:sec:0204}

在一个素理想处看不见某个元素，只说明有一个不属于该素理想的乘子零化它；若在每个极大理想处都看不见它，这些乘子生成的理想便不能包含于任何极大理想，只能是单位理想。由此可以检测元素是否为零、是否属于给定子模，以及给定同态是否单射、满射或同构。这些判据不要求模有限生成；把点处的有限生成元延伸到开邻域，则需要另行考察有限性。

\subsection{元素与模的消失判据}
\label{b3:subsec:0204:01}

\begin{theorem}\label{b3:thm:local-zero-detection}
设 \(R\) 为交换含幺环，\(M\) 为 \(R\) 模，\(m\in M\)。下列条件等价：
\begin{enumerate}
\item \(m=0\)；
\item 对每个素理想 \(\mathfrak p\)，\(m/1=0\) 在 \(M_{\mathfrak p}\) 中成立；
\item 对每个极大理想 \(\mathfrak m\)，\(m/1=0\) 在 \(M_{\mathfrak m}\) 中成立。
\end{enumerate}
因此 \(M=0\) 当且仅当所有 \(M_{\mathfrak m}\) 为零，也当且仅当所有 \(M_{\mathfrak p}\) 为零。
\end{theorem}
\begin{proof}
\((1)\Rightarrow(2)\Rightarrow(3)\) 直接成立，其中第二次蕴含用到极大理想是素理想。为证明 \((3)\Rightarrow(1)\)，设 \(m\ne0\)，则其零化子 \(\operatorname{Ann}_R(m)\) 是真理想，故包含于某个极大理想 \(\mathfrak m\)。若 \(m/1=0\) 在 \(M_{\mathfrak m}\) 中成立，\eqref{b3:eq:localized-element-zero}给出 \(s\notin\mathfrak m\) 且 \(sm=0\)，与 \(\operatorname{Ann}_R(m)\subseteq\mathfrak m\) 矛盾。因此第三个条件推出第一个。对 \(M\) 的每个元素应用该结论，即得模的断言。若 \(R\) 为零环，幺性模只能为零，空集上的检验也与结论相符。
\end{proof}

\begin{corollary}\label{b3:cor:submodule-local-detection}
设 \(R\) 为交换含幺环，\(M\) 为 \(R\) 模，\(N\subseteq M\) 为子模，\(m\in M\)。则
\[
m\in N\quad\Longleftrightarrow\quad
m/1\in N_{\mathfrak m}\text{ 对每个极大理想 }\mathfrak m.
\]
特别地，两个子模 \(N,L\subseteq M\) 相等，当且仅当它们在每个极大理想处局部化后相等。
\end{corollary}
\begin{proof}
由\cref{b3:thm:localization-exact}，\((M/N)_{\mathfrak m}\cong M_{\mathfrak m}/N_{\mathfrak m}\)。所以条件恰好表示 \(m+N\) 在每个局部化中为零；应用\cref{b3:thm:local-zero-detection}即得 \(m+N=0\)。对子模的每个元素分别应用包含判据，得到相等判据。
\end{proof}

例如整环 \(R\) 在分式域 \(K\) 内满足
\[
R=\bigcap_{\mathfrak m\text{ 极大}}R_{\mathfrak m}.
\]
确实，对 \(z\in K\)，分母理想 \(I_z=\{a\in R:az\in R\}\) 若为真理想，选包含它的极大理想 \(\mathfrak m\)。若 \(z\in R_{\mathfrak m}\)，则 \(z=b/s\)、\(s\notin\mathfrak m\)，于是 \(s\in I_z\)，矛盾。故处处属于局部环便迫使 \(I_z=R\)、\(z\in R\)。

\subsection{单射、满射与同构的局部检测}
\label{b3:subsec:0204:02}

\begin{theorem}\label{b3:thm:local-morphism-detection}
设 \(R\) 为交换含幺环，\(M,N\) 为 \(R\) 模。对任意 \(R\) 模同态 \(f:M\rightarrow N\)，单射、满射及同构这三种性质，均可在所有极大理想处的 \(f_{\mathfrak m}\) 上检测，也可在所有素理想处检测。更一般地，若 \(L,M,N\) 为 \(R\) 模，则复形
\[
L\xrightarrow{u}M\xrightarrow{v}N,\qquad vu=0,
\]
在 \(M\) 处正合，当且仅当每个极大理想处的局部化复形在中间项正合。
\end{theorem}
\begin{proof}
正合性定理给出
\[
(\ker f)_{\mathfrak m}\cong\ker f_{\mathfrak m},\qquad
(\operatorname{coker}f)_{\mathfrak m}\cong
\operatorname{coker}f_{\mathfrak m}.
\]
用\cref{b3:thm:local-zero-detection}分别检测这两个模是否为零，便得到单射与满射的判据；两者同时成立等价于同构。对复形同样考虑
\(H=(\ker v)/\operatorname{im}u\)。由局部化与核、像及商相容，其局部化恰为局部复形的中间同调。因此所有局部复形正合等价于所有 \(H_{\mathfrak m}=0\)，也就等价于 \(H=0\)。改用所有素理想的证明完全相同。
\end{proof}

上述结论检测的是一个已经给定的映射。若只知道每个 \(M_{\mathfrak m}\) 与 \(N_{\mathfrak m}\) 各自同构，却没有一个共同的全局映射，就不能直接应用该定理。局部同构之间是否相容，是另一项需要处理的数据。

\ProseParagraphBreak
给定一组元素 \(m_1,\ldots,m_r\in M\)，它们生成 \(M\)，当且仅当其像生成每个 \(M_{\mathfrak m}\)：把判据应用于 \(R^r\rightarrow M\)、\(e_i\mapsto m_i\) 的满射性即可。这里生成元是全局预先选定的；不能偷换成“每个点都能另选有限个局部生成元”。

\subsection{有限基本开覆盖与邻域有限性}
\label{b3:subsec:0204:03}

\begin{proposition}\label{b3:prop:principal-cover-finiteness}
设 \(R\) 为交换含幺环，\(M\) 为 \(R\) 模，\(f_1,\ldots,f_r\in R\)（\(r\ge1\)）生成单位理想。
\begin{enumerate}
\item 若 \(M_{f_i}=0\) 对每个 \(i\) 成立，则 \(M=0\)。
\item 若每个 \(M_{f_i}\) 都是有限生成 \(R_{f_i}\) 模，则 \(M\) 有限生成。
\end{enumerate}
\end{proposition}
\begin{proof}
对 \(m\in M\)，第一项的假设给出整数 \(n_i\ge1\) 使 \(f_i^{n_i}m=0\)。理想 \((f_1^{n_1},\ldots,f_r^{n_r})\) 仍为单位理想：否则它包含于某个极大理想，该极大理想由素性包含每个 \(f_i\)，与原假设矛盾。将 \(1\) 写成这些幂的线性组合，便得 \(m=0\)。

对第二项，把每个局部模的有限生成元写成 \(m_{ij}/f_i^{a_{ij}}\)。取由全部有限个分子 \(m_{ij}\) 生成的子模 \(N\)。由于 \(f_i\) 已可逆，\(N_{f_i}=M_{f_i}\)；正合性给出 \((M/N)_{f_i}=0\)。第一项应用于 \(M/N\)，即得 \(M=N\)。
\end{proof}

第一章中 \(D(f_i)\) 覆盖素谱恰好等价于 \((f_1,\ldots,f_r)=R\)。所以该命题说明，在一组有限基本开覆盖上给出有限生成性，足以得到全局有限生成性。若每个点有这样的基本开邻域，还可由素谱拟紧性抽取有限子覆盖。

\ProseParagraphBreak
仅要求每个茎 \(M_{\mathfrak p}\) 有限生成则不够。例如
\[
M=\bigoplus_{p\text{ 为素数}}\mathbb Z/p\mathbb Z
\]
不是有限生成整数模，因为有限个元素只涉及有限个直和分量；但 \(M_{(p)}\cong\mathbb Z/p\mathbb Z\)，而 \(M_{(0)}=0\)。另一方面，对每个非零整数 \(f\)，
\[
M_f\cong\bigoplus_{p\nmid f}\mathbb Z/p\mathbb Z.
\]
右边仍有无限多个非零分量，因而不是有限生成 \(\mathbb Z_f\) 模。因此所有茎都有限生成，并不能保证存在局部模有限生成的基本开邻域。这与\cref{b3:thm:support-finite}对有限生成性的要求相呼应。

\ProseParagraphBreak
此外，“在局部环上检验”不能随意改成“张量到剩余域上检验”。取域 \(k\)。在 \(R=k[t]_{(t)}\) 上，商映射 \(R/(t^2)\rightarrow R/(t)\) 不是同构，但张量到唯一的剩余域 \(k\) 后成为 \(k\) 上的恒等映射。下一章的 Nakayama 引理能在有限生成条件下用剩余向量空间检测生成性；它并不使剩余域张量变成一般的正合操作。

\subsection{相容局部元素的黏合}
\label{b3:subsec:0204:04}

前面的检测定理从已经给定的全局元素或同态出发。若先在有限个基本开集上给出模元素，还须检查它们在交集上相等，才能反过来构造全局元素。这个构造仍只需局部化的分式零判据与单位理想条件。

\begin{lemma}\label{b3:lem:localization-sheaf-gluing}
设 \(A\) 为交换含幺环，\(M\) 为 \(A\) 模，\(r\ge1\)，\(f_1,\ldots,f_r\in A\)，且 \(D(f_1),\ldots,D(f_r)\) 覆盖 \(\operatorname{Spec}A\)。在下列映射中，\(M\rightarrow\prod_{i=1}^rM_{f_i}\) 将 \(m\) 送到 \((m/1)_i\)；定义 \(\delta((s_i))_{ij}\) 时，先通过自然局部化映射将 \(s_i\in M_{f_i}\) 与 \(s_j\in M_{f_j}\) 送入 \(M_{f_if_j}\)，再取两者之差。则
\[
0\longrightarrow M\longrightarrow\prod_{i=1}^rM_{f_i}
\xrightarrow{\delta}\prod_{1\le i,j\le r}M_{f_if_j},
\qquad \delta((s_i))_{ij}=s_i-s_j
\]
在 \(M\) 及 \(\prod_{i=1}^rM_{f_i}\) 处正合。换言之，在各交集上相等的一组局部元素来自唯一的 \(M\) 中元素。
\end{lemma}
\begin{proof}
若 \(m\) 在各 \(M_{f_i}\) 中为零，\cref{b3:prop:principal-cover-finiteness}的第一项便给出 \(m=0\)，所以第一个映射单射。它的像在各交集上的限制相等，因而包含于 \(\ker\delta\)。

现设 \((s_i)\in\ker\delta\)。把每个 \(s_i\) 写成 \(u_i/f_i^{n_i}\)，其中 \(u_i\in M\)、\(n_i\ge0\)。因为指标 \(i\) 只有有限多个，可以选整数 \(N\ge1\)，使 \(N\ge n_i\) 对每个 \(i\) 成立。令 \(t_i=f_i^{N-n_i}u_i\)，则
\[
t_i/1=f_i^Ns_i\quad\text{在 }M_{f_i}\text{ 中成立}.
\]
在 \(M_{f_j}\) 中考虑 \(t_i/1-f_i^Ns_j\)。它进一步在 \(f_i\) 处局部化后为零，因为 \((M_{f_j})_{f_i}\cong M_{f_if_j}\)，且 \(s_i,s_j\) 在这一交集上的像相等。由分式零判据，对每对 \((i,j)\) 存在整数 \(L_{ij}\ge0\)，使 \(f_i^{L_{ij}}(t_i/1-f_i^Ns_j)=0\)。指标对也只有有限多个，故可选 \(L\ge L_{ij}\) 对所有 \((i,j)\) 成立。于是
\[
f_i^Lt_i/1=f_i^{N+L}s_j\quad\text{在 }M_{f_j}\text{ 中成立}.
\]
令 \(N'=N+L\)、\(t_i'=f_i^Lt_i\)，便得到 \(t_i'/1=f_i^{N'}s_j\) 在每个 \(M_{f_j}\) 中同时成立。此时 \((f_1^{N'},\ldots,f_r^{N'})=A\)：若这个理想包含于某个极大理想 \(\mathfrak m\)，由 \(N'\ge1\) 和素性可知所有 \(f_i\in\mathfrak m\)，这与 \(D(f_i)\) 覆盖全谱矛盾。因此可以选 \(a_1,\ldots,a_r\in A\)，使 \(\sum_i a_if_i^{N'}=1\)。取
\[
m=\sum_{i=1}^r a_it_i'\in M.
\]
对每个 \(j\)，在 \(M_{f_j}\) 中有
\[
m/1=\sum_{i=1}^r a_i(t_i'/1)
=\left(\sum_{i=1}^r a_if_i^{N'}\right)s_j=s_j.
\]
所以 \((s_i)\) 属于第一个映射的像，唯一性则由该映射的单射性得到。
\end{proof}

第二十四章构造仿射概形的结构层时，将用这条引理把基本开集上的相容截面拼成一个截面。

% !TeX root = ../../Book3.tex
% 纲要标题：### 2.5 局部环与局部同态
% 纲要位置：工作记录/计划纲要.md，第 1968 行。
% 纲要细目（写作范围）：
% * 局部环 \((R,\mathfrak m,\kappa)\) 的单位判据、极大理想与剩余域；局部性不蕴含 Noether 性
% * 局部同态及诱导的剩余域映射；有限型仿射代数的点处局部环
% * \(k[x_1,\ldots,x_n]_{(x_1,\ldots,x_n)}\)、离散赋值环及 Artin 局部环的对照
% * 先区分局部化、Hensel 化与完备化的任务；后续第13、23章分别建立相应理论
% ---

\section{局部环与局部同态}
\label{b3:sec:0205}

在 \(\mathfrak p\) 处局部化后，所有不属于 \(\mathfrak p\) 的元素都成为单位，留下的非单位恰好组成一个极大理想。这种情形值得独立命名，因为接下来的生成元判据和维数理论都将在这样的环上表述。局部性只说明单位与非单位的分界，不包含任何隐含的有限性假设。

\subsection{局部环的单位、极大理想与剩余域}
\label{b3:subsec:0205:01}

\begin{definition}\label{b3:def:local-ring}
设 \(R\) 为非零交换含幺环。若 \(R\) 恰有一个极大理想 \(\mathfrak m\)，则称 \(R\) 为\term[jubu-huan]{局部环}[local ring]。写作 \((R,\mathfrak m)\)，或连同剩余域 \(\kappa=R/\mathfrak m\) 写作 \((R,\mathfrak m,\kappa)\)。
\end{definition}
\symidx{\((R,\mathfrak m,\kappa)\)：局部环、极大理想及剩余域}

\begin{proposition}\label{b3:prop:local-ring-units}
设 \(R\) 为非零交换含幺环。下列条件等价：
\begin{enumerate}
\item \(R\) 为局部环；
\item 全部非单位组成一个理想；
\item 存在真理想 \(\mathfrak m\)，使 \(R\setminus\mathfrak m=R^\times\)。
\end{enumerate}
在这些条件下，该非单位理想就是唯一的极大理想；特别地，\(a\in\mathfrak m\) 时 \(1-a\) 是单位。
\end{proposition}
\begin{proof}
一个元素 \(a\) 非单位，当且仅当主理想 \((a)\) 为真理想，也就当且仅当 \(a\) 属于某个极大理想。若极大理想唯一，全部非单位恰为它，得到第二项与第三项。反过来，若非单位集是理想，则它包含每个真理想中的元素，又不含 \(1\)，所以是唯一的极大理想。第三项也给出同样结论。最后，若 \(a,1-a\) 都属于 \(\mathfrak m\)，则 \(1\in\mathfrak m\)，矛盾。
\end{proof}

局部环不必是 Noether 环。取域 \(k\)，令
\[
R=k[x_1,x_2,\ldots]_{(x_1,x_2,\ldots)}.
\]
这是由\cref{b3:prop:localized-prime-ideals}得到的局部环，但理想链
\((x_1)\subsetneq(x_1,x_2)\subsetneq\cdots\) 仍严格上升。若 \(x_{n+1}\) 属于前一项，清分母后有
\(s x_{n+1}\in(x_1,\ldots,x_n)\)，其中 \(s\) 的常数项非零；令前 \(n\) 个变量为零，在剩余变量的多项式整环中得到一个非零多项式乘 \(x_{n+1}\) 为零，矛盾。

\subsection{局部同态与点处局部环}
\label{b3:subsec:0205:02}

\begin{definition}\label{b3:def:local-homomorphism}
设 \((R,\mathfrak m)\) 与 \((A,\mathfrak n)\) 为局部环，\(\varphi:R\to A\) 为保单位环同态。若 \(\varphi(\mathfrak m)\subseteq\mathfrak n\)，则称 \(\varphi\) 为\term[jubu-tongtai]{局部同态}[local homomorphism]。
\end{definition}

任意含幺同态都把单位送到单位，所以总有
\(\varphi^{-1}(\mathfrak n)\subseteq\mathfrak m\)。因此局部同态的条件等价于
\(\varphi^{-1}(\mathfrak n)=\mathfrak m\)，也等价于非单位不会变成单位。它诱导一个含幺的剩余域同态
\[
R/\mathfrak m\longrightarrow A/\mathfrak n,
\]
该映射为单射，因为域到非零环的含幺同态核为零。

\begin{proposition}\label{b3:prop:induced-local-map}
设 \(R,A\) 为交换含幺环，\(\varphi:R\rightarrow A\) 为保单位环同态，\(\mathfrak q\in\operatorname{Spec}A\)，并令
\(\mathfrak p=\varphi^{-1}(\mathfrak q)\)，则有自然的局部同态
\[
R_{\mathfrak p}\longrightarrow A_{\mathfrak q},
\qquad r/s\longmapsto\varphi(r)/\varphi(s).
\]
它诱导剩余域的嵌入 \(\kappa(\mathfrak p)\rightarrow\kappa(\mathfrak q)\)。
\end{proposition}
\begin{proof}
\(s\notin\mathfrak p\) 意味着 \(\varphi(s)\notin\mathfrak q\)，因而目标分母合法；局部化的泛性质给出同态。若 \(r\in\mathfrak p\)，其像落在 \(\mathfrak qA_{\mathfrak q}\)，故该同态为局部同态。再取剩余域即得所述嵌入。
\end{proof}

例如，设 \(k\) 为域，\(I\) 为 \(k[x_1,\ldots,x_n]\) 的理想，则有限型仿射代数 \(A=k[x_1,\ldots,x_n]/I\) 在素理想 \(\mathfrak p\) 处的局部环为 \(A_{\mathfrak p}\)。对一个 \(k\) 有理点 \(a=(a_1,\ldots,a_n)\)，即所有 \(f\in I\) 在 \(a\) 处取零，代入映射的核 \(\mathfrak m_a\) 为极大理想，剩余域为 \(k\)。局部环中的元素可写成 \(f/g\)，要求 \(g(a)\ne0\)。这个描述表示在该点附近允许分母不消失，而不要求 \(g\) 在所有点上都非零。

\ProseParagraphBreak
局部环之间的任意含幺同态未必局部。例如 \(\mathbb Z_{(p)}\hookrightarrow\mathbb Q\) 把非单位 \(p\) 变为单位，其极大理想不能映入 \(\mathbb Q\) 的极大理想 \(0\)。所以使用剩余域映射前必须检查局部条件。

\subsection{多项式局部环、离散赋值环与 Artin 局部环}
\label{b3:subsec:0205:03}

设 \(k\) 为域，\(n\ge1\)。在
\[
R=k[x_1,\ldots,x_n]_{(x_1,\ldots,x_n)}
\]
中，唯一极大理想由 \(x_1,\ldots,x_n\) 生成，剩余域为 \(k\)。当 \(n=1\) 时，每个非零元素都唯一写成 \(x^r u\)，其中 \(r\ge0\)、\(u\) 为单位。因为多项式分子可提取唯一的 \(x\) 次幂，而局部化的分母常数项非零。非零理想中取这样的最小 \(r\)，便知该理想等于 \((x^r)\)。

这里的指数 \(r\) 测量函数在原点的零点阶数。将它推广到分式域 \(k(x)\) 时，若非零多项式 \(f=x^a f_0\)、\(g=x^b g_0\)，且 \(f_0(0),g_0(0)\ne0\)，就令
\[
v(f/g)=a-b.
\]
这是分子与分母零点阶数之差；若另一个分式表示同一元素，交叉相乘使这两个差相等，所以它不依赖表示。乘积使阶数相加，非零和的阶数则至少为两项阶数的较小者，因为通分后最低次项可能相消。于是 \(v(x^r u)=r\)，局部环正好由阶数非负的有理函数及零组成，而其中的单位恰好阶数为零。离散赋值环将这种整数阶数及其运算规则抽象出来。

\begin{definition}\label{b3:def:discrete-valuation-example}
设 \(K\) 为域，满射 \(v:K^\times\rightarrow\mathbb Z\) 满足
\[
\begin{aligned}
v(ab)&=v(a)+v(b) &&(a,b\in K^\times),\\
v(a+b)&\ge\min\bigl\{v(a),v(b)\bigr\} &&(a,b\in K^\times,\ a+b\ne0).
\end{aligned}
\]
这样的 \(v\) 称为一个\term[lisan-fuzhi]{离散赋值}[discrete valuation]；相应子环
\(V=\{0\}\cup\bigl\{\,a\in K^\times\mid v(a)\ge0\,\bigr\}\) 称为\term[lisan-fuzhi-huan]{离散赋值环}[discrete valuation ring]，简称 DVR。
\end{definition}

由赋值条件，\(V\) 确实对加法及乘法封闭；其单位恰为赋值为零的元素。若选 \(\pi\) 满足 \(v(\pi)=1\)，每个非零元素为 \(\pi^r u\)，并且非单位组成 \((\pi)\)。因此它是非域的局部整环，所有非零理想均为 \((\pi^r)\)，证明与上一段的最小赋值论证相同。取素数 \(p\)，则 \(\mathbb Z_{(p)}\) 对应有理数的 \(p\) 指数；\(k[x]_{(x)}\) 对应有理函数在 \(x=0\) 处的零点阶数。DVR 的其他内在刻画见\chapref{b3:ch:10}。

\ProseParagraphBreak
另一类局部环来自有限维代数。仍设 \(k\) 为域，对整数 \(r\ge2\)，取
\[
B=k[\epsilon]/(\epsilon^r),\qquad
C=k[x,y]/(x,y)^2.
\]
在 \(B\) 中，一个元素为单位当且仅当常数项非零：把它写成 \(a(1+n)\)，其中 \(a\in k^\times\)、\(n^r=0\)，逆为
\(a^{-1}\bigl(1-n+\cdots+(-n)^{r-1}\bigr)\)。唯一极大理想 \((\bar\epsilon)\) 的 \(r\) 次幂为零，而 \(\bar\epsilon\ne0\)，所以 \(B\) 不是整环；若取 \(r=1\)，则商环就是域 \(k\)。在 \(C\) 中同理，唯一极大理想 \(\mathfrak m=(\bar x,\bar y)\) 的平方为零。极大理想 \(\mathfrak m\) 作为 \(k\) 向量空间以 \(\bar x,\bar y\) 为基，而整个环 \(C\) 以 \(1,\bar x,\bar y\) 为基。\(C\) 对 \(\mathfrak m\) 的作用只通过常数项，故 \(\mathfrak m\) 中任何一个元素所生成的理想至多一维；这个极大理想必须用两个元素生成。这些环作为 \(k\) 向量空间有限维，故理想降链必稳定，都是 Artin 局部环。

\ProseParagraphBreak
\cref{b3:tab:local-ring-examples}列出三种情形。极大理想只有一个，并不决定该理想是否主生成、是否幂零，也不决定环是否为整环。DVR 的极大理想由一个非零元生成，但它的任何次幂都非零；后两类环则具有非零幂零元。

\begin{table}[htbp]
\centering
\caption{局部环的三个例子}
\label{b3:tab:local-ring-examples}
\begin{tabularx}{\textwidth}{@{}>{\raggedright\arraybackslash}p{.30\textwidth}>{\raggedright\arraybackslash}p{.23\textwidth}*{3}{>{\raggedright\arraybackslash}X}@{}}
\toprule
环 & 极大理想 \(\mathfrak m\) & 整环 & \(\mathfrak m\) 主生成 & \(\mathfrak m\) 幂零 \\
\midrule
\(k[x]_{(x)}\) & \((x)\) & 是 & 是 & 否 \\
\(k[\epsilon]/(\epsilon^r)\)，\(r\ge2\) & \((\bar\epsilon)\) & 否 & 是 & 是 \\
\(k[x,y]/(x,y)^2\) & \((\bar x,\bar y)\) & 否 & 否 & 是 \\
\bottomrule
\end{tabularx}
\end{table}

\subsection{局部化、Hensel 化与完备化的比较}
\label{b3:subsec:0205:04}

设 \(k\) 为域，\(R=k[x]_{(x)}\)。在 \(x=0\) 附近作分式计算，只需允许常数项非零的多项式作分母。例如
\(\dfrac{1}{1-x}\in R\)，但一个任意给定的形式幂级数未必能写成这样的有理函数。局部化保留的是分式计算，而不是所有无穷次近似。

\ProseParagraphBreak
若希望同时保存模掉 \(x,x^2,\ldots\) 的相容近似，就要考虑序列 \((a_n)_{n\ge1}\)，其中 \(a_n\in R/(x^n)\)，且 \(a_{n+1}\) 模 \(x^n\) 的像为 \(a_n\)。逐项相加、相乘使这些序列成为一个环。对一般局部环 \((A,\mathfrak m)\)，以同样方式由 \(A/\mathfrak m^n\) 构成的环记为
\[
\widehat A=\varprojlim_{n\ge1} A/\mathfrak m^n,
\]
称为 \(A\) 的\term[wanbeihua]{完备化}[completion]。对上面的 \(R\)，每个商同构于 \(k[x]/(x^n)\)：常数项非零的多项式在这个商中已经可逆。相容的截断多项式逐项确定一个形式幂级数，反过来每个级数也给出这样一组截断，因此 \(\widehat R\cong k[\![x]\!]\)。例如，\(\dfrac{1}{1-x}\) 对应各阶截断 \(1+x+\cdots+x^{n-1}\)，也就是级数 \(\sum_{j\ge0}x^j\)。完备化与正合性、原环到完备化的映射等问题见\chapref{b3:ch:13}。

\ProseParagraphBreak
还有一种要求来自解方程。再设 \(\operatorname{char}k\ne2\)，则方程 \(T^2-(1+x)=0\) 模掉 \(x\) 后在 \(T=1\) 处有单根；但它在 \(R\) 中没有根，因为有理函数的平方在不可约多项式 \(1+x\) 处具有偶数阶，而 \(1+x\) 的阶为一。在 \(k[\![x]\!]\) 中，若要求根的常数项为 \(1\)，比较各次系数时新系数总乘以 \(2\)，便可逐次唯一确定这个根。

\ProseParagraphBreak
第23章的\cref{b3:thm:henselization-existence}将给出 \(R\) 的 Hensel 化 \(R^{\mathrm h}\)。它保证剩余域中的单根或互素因子能够提升，因而其中也有上述方程的根，且其剩余类为 \(1\)。两种扩张都能解决这个方程，但所加入的数据不同：完备化接纳所有相容的有限阶近似，Hensel 化则以\cref{b3:def:henselization}的泛性质规定解决这类局部代数提升问题所需的扩张，不要求接纳任意给定的形式幂级数。这也是在完备化之外另行构造 Hensel 化的原因；其与完备化的关系将在第23章具体说明。


\begin{chaptersummary}
局部化中的等式允许再乘一个分母后成立；这一零判据同时解释了分式模的构造、局部化的正合性和理想收缩时的饱和。从局部模 \(M_{\mathfrak p}\) 取剩余向量空间，还要模去 \(\mathfrak pM_{\mathfrak p}\)，得到纤维 \(M\otimes_R\kappa(\mathfrak p)\)。这次取商会丢失局部模中的信息，也未必保持单射。

\ProseParagraphBreak
对任意模，元素是否为零、是否属于给定子模、两个子模是否相等，以及给定同态是否单射、满射或同构，都能在所有极大理想处检测；给定复形的正合性也如此。这些判据不要求有限生成。有限生成使支集等于零化子的闭集，并使点处的一组全局生成元能在某个基本开邻域上继续生成；有限展示进一步使有限个关系可以统一清分母，从而得到 Hom 的局部化同构。仅有每个局部模有限生成，仍不足以保证邻域上的有限生成性。

\ProseParagraphBreak
有限基本开覆盖上的局部模元素若在两两交集上一致，清除有限多个分母后就能拼合成唯一的全局元素。这里先构造的是元素；从各处分别给出的抽象同构构造全局对象，还须处理同构之间的相容性。

\ProseParagraphBreak
局部环把非单位集中在唯一极大理想中，为下一章的 Nakayama 引理准备了单位判据。多项式局部环、DVR 及 Artin 局部环已经展示不同的极大理想结构；完备化与 Hensel 化各自增加的内容将在后续章节另行建立。
\end{chaptersummary}

\BookExercises

\begin{exercise}\label{b3:ex:02-cyclic-integers}
设 \(n\ge1\)，\(p\) 为素数，写 \(n=p^a d\)、\(p\nmid d\)。求
\((\mathbb Z/n\mathbb Z)_{(p)}\)，并说明 \(a=0\) 时的结果。
\end{exercise}
\begin{solution}
在 \(\mathbb Z_{(p)}\) 中 \(d\) 是单位，所以
\(n\mathbb Z_{(p)}=p^a\mathbb Z_{(p)}\)。由商与局部化相容，
\[
(\mathbb Z/n\mathbb Z)_{(p)}
\cong\mathbb Z_{(p)}/p^a\mathbb Z_{(p)}
\cong\mathbb Z/p^a\mathbb Z.
\]
末个同构把 \(b/s\) 送到 \(b s^{-1}\) 模 \(p^a\) 的类，\(p\nmid s\) 保证逆存在；核恰为 \(p^a\mathbb Z_{(p)}\)。若 \(a=0\)，\(n\) 已是单位，商为零。
\end{solution}

\begin{exercise}\label{b3:ex:02-fraction-kernel}
设 \(k\) 为域，令 \(R=k[x,y]/(xy)\)，仍用 \(x,y\) 记剩余类。求 \(R\rightarrow R_x\) 的核，并显式确定 \(R_x\)。
\end{exercise}
\begin{solution}
\(x\) 可逆后，由 \(xy=0\) 得 \(y=0\)，所以 \(R_x\cong k[x,x^{-1}]\)。每个 \(R\) 中的元素唯一写成
\(a(x)+y\cdot b(y)\)。其在局部化中的像为 \(a(x)\)，由 \(k[x]\) 到 Laurent 多项式环的嵌入可知像为零当且仅当 \(a=0\)。故核为 \((y)\)，恰好是被某个 \(x\) 幂零化的元素组成的理想。
\end{solution}

\begin{exercise}\label{b3:ex:02-infinite-product}
固定素数 \(p\)，取 \(M=\prod_{n\ge1}\mathbb Z/p^n\mathbb Z\)，把 \(p\) 变为单位。证明每个分量的局部化均为零，而 \(M[1/p]\ne0\)。这与局部化正合有无矛盾？
\end{exercise}
\begin{solution}
第 \(n\) 个分量被 \(p^n\) 零化，故局部化为零。取 \(m=(1\bmod p^n)_n\)。对任何 \(r\)，其第 \(r+1\) 个分量满足
\(p^r\not\equiv0\pmod{p^{r+1}}\)，所以 \(p^rm\ne0\)。由零判据，\(m/1\ne0\)。
因此局部化不必与无限直积交换。正合性讨论的是核、像和有限的正合条件，不包含任意无限直积相容性，故无矛盾。
\end{solution}

\begin{exercise}\label{b3:ex:02-saturation-computation}
设 \(k\) 为域。在 \(R=k[x,y]\) 中，令 \(I=(x^2,xy)\)。求 \(I:y^\infty\)、\(I:x^\infty\)，并写出在 \(R_y,R_x\) 中的扩张理想。
\end{exercise}
\begin{solution}
因 \(yx\in I\)，有 \((x)\subseteq I:y^\infty\)。若 \(y^n f\in I\subseteq(x)\)，模掉 \(x\) 后在整环 \(k[y]\) 中得到 \(y^n\bar f=0\)，故 \(\bar f=0\)，即 \(f\in(x)\)。所以 \(I:y^\infty=(x)\)，相应扩张为 \((x)R_y\)。
又 \(x^2\in I\)，故 \(1\in I:x^\infty\)，得到 \(I:x^\infty=R\)、\(IR_x=R_x\)。
\end{solution}

\begin{exercise}\label{b3:ex:02-two-saturations}
设 \(k\) 为域，在 \(k[x,y]\) 中取 \(I=(xy)\)、\(J=(x,y)\)，并令 \(S\) 为 \(x,y\) 生成的乘法集。比较 \(I:J^\infty\) 与 \(\operatorname{Sat}_S(I)\)。
\end{exercise}
\begin{solution}
若 \(fJ^n\subseteq(xy)\)，则 \(fx^n\) 和 \(fy^n\) 都被 \(xy\) 整除。第一式模掉 \(y\) 后迫使 \(f\in(y)\)，第二式模掉 \(x\) 后迫使 \(f\in(x)\)。多项式的单项式展开给出
\((x)\cap(y)=(xy)\)，故 \(f\in I\)。反向显然，因而 \(I:J^\infty=I\)。
另一方面 \(xy\in I\cap S\)，所以 \(\operatorname{Sat}_S(I)=R\)。前者要求一整个幂理想都送入 \(I\)，后者只要求找到一个乘子。
\end{solution}

\begin{exercise}\label{b3:ex:02-primes-invert-six}
列出 \(\mathbb Z[1/6]\) 的素理想，并求
\((\mathbb Z/42\mathbb Z)\otimes_{\mathbb Z}\mathbb Z[1/6]\) 的支集。
\end{exercise}
\begin{solution}
素理想对应于 \(\mathbb Z\) 中与 \(\{6^n:n\ge0\}\) 不交的素理想，故为
\(0\) 及 \((p)\mathbb Z[1/6]\)，其中 \(p\ne2,3\) 为素数。局部化后 \(42=6\cdot7\) 与 \(7\) 生成同一理想，所以该模为 \(\mathbb Z[1/6]/(7)\cong\mathbb F_7\)。其支集只有极大理想 \((7)\mathbb Z[1/6]\)。
\end{solution}

\begin{exercise}\label{b3:ex:02-residue-generic}
设 \(k\) 为域。在 \(R=k[x,y]\) 中取 \(\mathfrak p=(x)\)、\(\mathfrak m=(x,y)\)。分别求剩余域，说明为什么 \(R_{\mathfrak p}\) 中 \(y\) 可逆，而 \(R_{\mathfrak m}\) 中不可逆。
\end{exercise}
\begin{solution}
\(R/\mathfrak p=k[y]\)，故 \(\kappa(\mathfrak p)=k(y)\)；\(R/\mathfrak m=k\)，故
\(\kappa(\mathfrak m)=k\)。因 \(y\notin\mathfrak p\)，它是 \(R_{\mathfrak p}\) 的允许分母；而 \(y\in\mathfrak m\)，在 \(R_{\mathfrak m}\) 中属于唯一极大理想，不能为单位。素理想处的剩余域不必是基域，也不必由一个闭点给出。
\end{solution}

\begin{exercise}\label{b3:ex:02-hyperbola-fibres}
设 \(k\) 为域，\(\varphi:k[t]\rightarrow k[x,y]\) 为满足 \(\varphi(t)=xy\) 的 \(k\)-代数同态。求 \(\varphi\) 在素理想 \((t-a)\)、\(a\in k\)，处的纤维代数，并分别讨论 \(a=0\) 与 \(a\ne0\)。
\end{exercise}
\begin{solution}
纤维代数为 \(k[x,y]/(xy-a)\)。若 \(a=0\)，素谱是 \(V(x)\cup V(y)\)，且
\(x,y\) 的非零剩余类乘积为零。若 \(a\ne0\)，关系说明 \(x\) 可逆，逆为 \(a^{-1}y\)；映射
\[
k[x,y]/(xy-a)\longrightarrow k[x,x^{-1}],\qquad
x\longmapsto x,\quad y\longmapsto ax^{-1}
\]
的逆把 \(x^{-1}\) 送到 \(a^{-1}y\)，故为同构。此时纤维环是整环，且不包含 \(x=0\) 的点。
\end{solution}

\begin{exercise}\label{b3:ex:02-prufer-support}
设 \(p\) 为素数。求整数模 \(\mathbb Z[1/p]/\mathbb Z\) 的零化子与支集，并据此说明无限生成模的支集与零化子闭集可以不同。
\end{exercise}
\begin{solution}
没有非零整数能同时零化全部 \(p^{-n}+\mathbb Z\)，所以零化子为零。在 \((q)\)、\(q\ne p\)，处，\(p\) 已可逆，每个元素被某个 \(p\) 幂零化，故局部模为零；在 \((0)\) 处同样为零。在 \((p)\) 处，任何不被 \(p\) 整除的整数作用在每个有限循环子群上都是可逆的，因此自然局部化不消去非零元素。支集恰为 \(\bigl\{(p)\bigr\}\)，而 \(V(0)=\operatorname{Spec}\mathbb Z\)。
\end{solution}

\begin{exercise}\label{b3:ex:02-zero-neighbourhood}
设 \(R\) 为交换含幺环，\(M\) 为有限生成 \(R\) 模，\(\mathfrak p\) 为 \(R\) 的素理想。若 \(M_{\mathfrak p}=0\)，证明存在 \(f\notin\mathfrak p\) 使 \(M_f=0\)。更一般地，若 \(m_1,\ldots,m_r\in M\) 的像生成 \(M_{\mathfrak p}\)，证明存在 \(f\notin\mathfrak p\)，使 \(m_1/1,\ldots,m_r/1\) 生成 \(M_f\)。
\end{exercise}
\begin{solution}
由\cref{b3:thm:support-finite}，存在 \(f\notin\mathfrak p\) 使 \(fM=0\)，因而 \(M_f=0\)。对第二问，令 \(N\) 为这些全局元素生成的子模。有限生成商模 \(Q=M/N\) 满足 \(Q_{\mathfrak p}=0\)，故存在 \(f\notin\mathfrak p\) 使 \(Q_f=0\)。由局部化正合性，\(N_f=M_f\)，正是所需结论。
\end{solution}

\begin{exercise}\label{b3:ex:02-hom-finite-generated-insufficient}
设 \(k\) 为域，令 \(R=\prod_{n\ge1}k\)，\(I=\bigoplus_{n\ge1}k\subseteq R\)。对有限集 \(F\subseteq\mathbb Z_{\ge1}\)，令 \(e_F\) 在 \(F\) 的坐标上等于 \(1\)，其余坐标上等于 \(0\)，并取 \(S=\{1-e_F:F\;\text{有限}\}\)。再令 \(e\in R\) 在偶数坐标上等于 \(1\)，奇数坐标上等于 \(0\)，\(J=eI\)、\(M=R/I\)、\(N=R/J\)。计算自然映射
\[
S^{-1}\operatorname{Hom}_R(M,N)\longrightarrow
\operatorname{Hom}_{S^{-1}R}(S^{-1}M,S^{-1}N)
\]
的像与余核，说明本题中源与像都非零，但映射仍不满射。
\end{exercise}
\begin{solution}
记 \(A=S^{-1}R\)。由\cref{b3:prop:hom-localization-fg-counterexample}中的同构，\(A\cong R/I\)，自然映射把 \(a/s\) 送到 \(a+I\)。由于 \(J\subseteq I\)，有 \(S^{-1}J=0\)；局部化商模的正合列给出 \(S^{-1}M\cong A\)、\(S^{-1}N\cong A\)。因此目标同态模由同态在 \(1\) 上的值识别为 \(A\)。

一个同态 \(M\to N\) 由 \(1+I\) 的像 \(a+J\) 决定，存在的条件是 \(Ia\subseteq J\)。用各个单坐标幂等元相乘可知，这等价于 \(a\) 的所有奇数坐标均为零：偶数坐标产生的乘积已属于 \(J\)，奇数坐标产生的乘积只能为零。所以
\[
\operatorname{Hom}_R(M,N)\cong eR/J.
\]
局部化后，\(S^{-1}(eR/J)\cong eA\)。自然映射把在 \(1+I\) 上取值为 \(a+J\) 的同态送到在 \(1\) 上取值为 \(a/1\) 的同态；每个 \(s\in S\) 在 \(A\) 中等于 \(1\)，所以在上述识别下它就是包含 \(eA\hookrightarrow A\)。其像为 \(eA\)，余核为
\[
A/eA\cong(1-e)A.
\]
最后，\(e\) 与 \(1-e\) 都有无限多个非零坐标，均不属于 \(I\)，所以它们在 \(A\cong R/I\) 中的像均非零。因而源与像非零，余核也非零。例如，\(A\to A\)、\(a\mapsto(1-e)a\) 就不能由局部化前的同态统一清分母得到。
\end{solution}

\begin{exercise}\label{b3:ex:02-unbounded-hom-denominators}
取 \(M=\bigoplus_{n\ge1}\mathbb Z e_n\)、\(N=\mathbb Z\)，把一个素数 \(p\) 变为单位。证明局部化后由 \(e_n\mapsto p^{-n}\) 定义的同态，不来自
\(\operatorname{Hom}_{\mathbb Z}(M,\mathbb Z)[1/p]\)。
\end{exercise}
\begin{solution}
局部化后 \(M[1/p]\) 是具有基 \(e_n\) 的自由 \(\mathbb Z[1/p]\) 模，因此指定这些像确实定义同态。若它来自 \(f/p^r\)，则对每个 \(n\) 有
\(f(e_n)=p^{r-n}\) 作为有理数成立。但 \(n>r\) 时右侧不是整数，矛盾。这个例子的困难已经发生在无限个生成元上，无需再考虑生成关系。
\end{solution}

\begin{exercise}\label{b3:ex:02-intersection-calculation}
设 \(k\) 为域。在 \(k[x,y]\) 中，先求 \((x)\cap(y)\)，再在 \(y\) 处局部化，核对结果等于两个扩张理想的交。另取素数 \(p\)，比较整数环中 \(\bigcap_{n\ge1}(p^n)\) 的局部化与各个 \((p^n)\) 在 \(\mathbb Z[1/p]\) 中的扩张理想之交。
\end{exercise}
\begin{solution}
属于 \((x)\) 的多项式的每个非零单项式都含 \(x\)，属于 \((y)\) 的则都含 \(y\)。因此交为 \((xy)\)。在 \(R_y\) 中，\((y)R_y=R_y\)，而
\((xy)R_y=(x)R_y\)，故
\[
\bigl((x)\cap(y)\bigr)R_y=(x)R_y=(x)R_y\cap(y)R_y.
\]
这与有限交的局部化相容性一致。对于第二问，\(\bigcap_{n\ge1}p^n\mathbb Z=0\)，因为一个非零整数不能被任意高次的 \(p\) 幂整除；但 \(p\) 可逆后，每个 \(p^n\mathbb Z[1/p]=\mathbb Z[1/p]\)。于是
\[
\left(\bigcap_{n\ge1}p^n\mathbb Z\right)[1/p]=0,
\qquad
\bigcap_{n\ge1}p^n\mathbb Z[1/p]=\mathbb Z[1/p]\ne0.
\]
因此局部化与有限交相容，却不必与无限交相容。
\end{solution}

\begin{exercise}\label{b3:ex:02-omit-prime}
固定素数 \(q\)。在 \(\mathbb Q\) 内计算 \(\bigcap_{p\ne q}\mathbb Z_{(p)}\)，其中 \(p\) 遍历其余素数。
\end{exercise}
\begin{solution}
把一个有理数写成互素分子分母 \(a/b\)、\(b>0\)。它属于 \(\mathbb Z_{(p)}\) 当且仅当 \(p\nmid b\)。要求对所有 \(p\ne q\) 成立，就等价于 \(b\) 只有素因子 \(q\)。因此交恰为 \(\mathbb Z[1/q]\)。删去一个极大理想处的检测，确实可能失去信息。
\end{solution}

\begin{exercise}\label{b3:ex:02-two-open-membership}
设 \(R\) 为交换含幺环，\(M\) 为 \(R\) 模，\(N\subseteq M\) 为子模，\(m\in M\)，且 \(f,g\in R\) 满足 \(fR+gR=R\)。若 \(m/1\in N_f\) 且 \(m/1\in N_g\)，直接用分母证明 \(m\in N\)，不要只引用全部极大理想处的判据。
\end{exercise}
\begin{solution}
在 \(M/N\) 中，\(m+N\) 分别在 \(f,g\) 处消失，所以存在 \(r,s\ge1\) 使
\(f^rm,g^sm\in N\)。理想 \((f^r,g^s)\) 为单位理想：否则一个包含它的极大理想将同时包含 \(f,g\)。取 \(af^r+bg^s=1\)，得到
\(m=a(f^rm)+b(g^sm)\in N\)。
\end{solution}

\begin{exercise}\label{b3:ex:02-stalk-versus-neighbourhood}
令 \(M=\bigoplus_p\mathbb Z/p\mathbb Z\)。证明对每个非零整数 \(n\)，\(M[1/n]\) 仍不是有限生成 \(\mathbb Z[1/n]\) 模。由此解释为什么 \(M_{(0)}=0\) 不能在 \((0)\) 的一个基本开邻域上给出有限生成性。
\end{exercise}
\begin{solution}
若 \(p\mid n\)，第 \(p\) 个分量被局部化消去；若 \(p\nmid n\)，\(n\) 在
\(\mathbb Z/p\mathbb Z\) 上已可逆，该分量保留。故
\(M[1/n]\cong\bigoplus_{p\nmid n}\mathbb Z/p\mathbb Z\)。
仍有无限多个非零分量，而有限个元素的 \(\mathbb Z[1/n]\) 线性组合只能涉及有限个分量，所以不有限生成。包含 \((0)\) 的基本开集恰好是非零 \(n\) 给出的 \(D(n)\)，因此没有一个这样的邻域满足所需有限性。
\end{solution}

\begin{exercise}\label{b3:ex:02-fibre-misses-kernel}
设 \(k\) 为域，\(R=k[t]_{(t)}\)。计算 \(R/(t^3)\rightarrow R/(t)\) 的核，并证明张量到剩余域后该核的自然像为零，而张量后的商映射是同构。
\end{exercise}
\begin{solution}
核为 \((t)/(t^3)\)，其中 \(t+(t^3)\ne0\)。张量到 \(k=R/(t)\) 时，两个商模都变成 \(k\)，诱导映射为恒等。核的张量为
\[
\bigl((t)/(t^3)\bigr)\otimes_R k
\cong(t)/(t^2)\cong k,
\]
但它到 \(\bigl(R/(t^3)\bigr)\otimes_R k=k\) 的自然映射把 \(t\) 的类送到零。于是原序列的单射在张量后不再单射；这恰是剩余域张量不能替代局部化正合性的原因。
\end{solution}

\begin{exercise}\label{b3:ex:02-local-map-examples}
设 \(k\) 为域。证明 \(k[x]_{(x)}\rightarrow k[x,y]_{(x,y)}\) 的自然同态为局部同态，并求其剩余域映射。再说明 \(k[x]_{(x)}\hookrightarrow k(x)\) 为什么不是局部同态。
\end{exercise}
\begin{solution}
来源中允许的分母在 \(x=0\) 处非零，进入目标后在 \((0,0)\) 处仍非零，所以同态存在。极大理想 \((x)\) 映入 \((x,y)\)，故同态局部；两边剩余域都是 \(k\)，诱导映射为恒等。到分式域的嵌入把非单位 \(x\) 变为单位，或者说 \((x)\) 的像不包含于目标极大理想 \(0\)，所以不局部。
\end{solution}

\begin{exercise}\label{b3:ex:02-dvr-ideals}
设 \(k\) 为域。在 \(R=k[x]_{(x)}\) 中求理想 \(\bigl(x^3(1+x),x^5\bigr)\)，并求商环的所有理想及其作为 \(k\) 向量空间的维数。
\end{exercise}
\begin{solution}
\(1+x\) 为单位，所以原理想等于 \((x^3)\)。\(R\) 的每个非零理想为一个
\((x^r)\)，因此商环 \(R/(x^3)\) 的理想为
\(0,(\bar x^2),(\bar x),R/(x^3)\)。该商同构于 \(k[x]/(x^3)\)，基为
\(1,\bar x,\bar x^2\)，维数为 \(3\)；其中三个非零理想的维数分别为 \(1,2,3\)。
\end{solution}

\begin{exercise}\label{b3:ex:02-square-zero-local}
设 \(k\) 为域。对 \(A=k[x,y]/(x,y)^2\)，写出一般单位的逆，并证明其极大理想不能由一个元素生成。
\end{exercise}
\begin{solution}
一般元素为 \(a+b\bar x+c\bar y\)。若 \(a\ne0\)，令 \(n=b\bar x+c\bar y\)，有
\(n^2=0\)，故逆为 \(a^{-1}-a^{-2}n\)。若 \(a=0\)，该元素平方为零，不能为单位。所以极大理想为 \(\mathfrak m=k\bar x\oplus k\bar y\)。因 \(\mathfrak m^2=0\)，一个元素 \(u\in\mathfrak m\) 生成的理想只有 \(ku\)，至多一维，不能等于二维的 \(\mathfrak m\)。
\end{solution}

\begin{exercise}\label{b3:ex:02-completion-compatible}
设 \(k\) 为域。直接从相容序列定义构造
\(\varprojlim_{n\ge1} k[x]/(x^n)\cong k[\![x]\!]\)，并说明 \(1/(1-x)\) 在该同构下的像。再证明自然映射 \(k[x]_{(x)}\rightarrow k[\![x]\!]\) 单射但不满射，构造一个不在其像中的级数。
\end{exercise}
\begin{solution}
一个相容序列可唯一写成
\(a_0+a_1x+\cdots+a_{n-1}x^{n-1}\pmod{x^n}\)，因为相邻项的低次系数必须相同。把它送到 \(\sum_{j\ge0}a_jx^j\) 得到双射；每个固定次数的加法与乘法只涉及有限个系数，故双射保留环运算。
在每个商环中，\(1-x\) 的逆为 \(1+x+\cdots+x^{n-1}\)，因此相容序列对应
\(\sum_{j\ge0}x^j\)。这是形式系数的等式，不涉及实数或复数意义的收敛。

常数项非零的多项式在 \(k[\![x]\!]\) 中可逆，因此把 \(f/g\in k[x]_{(x)}\) 送到 \(f\) 与 \(g\) 的形式逆的乘积，就得到所述自然映射。若其像为零，乘以 \(g\) 得 \(f=0\)，故映射单射。取
\[
h(x)=\sum_{j\ge1}x^{j!}.
\]
若 \(h=f/g\)，其中 \(f,g\in k[x]\)、\(g(0)\ne0\)，则 \(gh=f\)。选充分大的 \(j\)，使 \(j!>\deg f\) 且 \(j!-(j-1)!>\deg g\)。在 \(h\) 中，\(x^{j!}\) 前面的 \(\deg g\) 个系数全为零，故 \(gh\) 的 \(x^{j!}\) 系数等于 \(g(0)\ne0\)，而 \(f\) 的该系数为零，矛盾。所以 \(h\) 不在像中，映射不满射。
\end{solution}

\begin{exercise}\label{b3:ex:02-root-lifting-warning}
设 \(k\) 为特征不为 \(2\) 的域。证明 \(T^2-(1+x)\) 在 \(k[\![x]\!]\) 中有唯一常数项为 \(1\) 的根，却在 \(k(x)\) 中没有根。
\end{exercise}
\begin{solution}
令根为 \(1+\sum_{n\ge1}a_nx^n\)。比较一次项得 \(2a_1=1\)；对 \(n\ge2\) 比较系数得
\[
2a_n+\sum_{i=1}^{n-1}a_i a_{n-i}=0.
\]
因 \(2\) 可逆，这些等式递归且唯一地确定全部系数，并使平方恰为 \(1+x\)。
若有有理函数根 \(f/g\)，约成互素多项式后由 \(f^2=(1+x)g^2\)，不可约多项式 \(1+x\) 在左侧的指数为偶数，在右侧为奇数，矛盾。因此完备化中可以出现原局部环及其分式域都没有的代数根。
\end{solution}

\begin{exercise}\label{b3:ex:02-two-open-gluing-calculation}
令 \(R=\mathbb Z/12\mathbb Z\)，\(f=\bar2\)、\(g=\bar3\)。在 \(R_f\) 和 \(R_g\) 中分别给出 \(s_f=1/f\)、\(s_g=1/g\)。证明这两个元素在 \(R_{fg}\) 中相容，并直接清分母构造其唯一的全局元素。具体要求找出共同指数 \(N\)、元素 \(u,v\in R\) 和系数 \(a,b\in R\)，使
\[
u/1=f^Ns_f,\qquad v/1=g^Ns_g,
\qquad af^N+bg^N=1,
\]
且 \(m=au+bv\) 的两个局部像分别为 \(s_f,s_g\)；不能仅引用黏合引理或中国剩余定理。
\end{exercise}
\begin{solution}
由 \(-f+g=1\)，\(D(f),D(g)\) 覆盖全谱。又 \((fg)^2=\bar{36}=0\)，故 \(R_{fg}=0\)，两个局部元素在交叠上相等。

取 \(N=3\)、\(u=f^2=\bar4\)、\(v=g^2=\bar9\)，便有所要求的两个分母等式。整数恒等式 \(-10\cdot8+3\cdot27=1\) 给出 \(a=\overline{-10}\)、\(b=\bar3\)。在 \(R_f\) 中，由 \(f^2g=0\) 且 \(f\) 可逆，得 \(g=0\)，所以 \(v=0\)；把系数等式除以 \(f\)，得到 \(m/1=af^2=1/f\)。在 \(R_g\) 中，由 \(f^2g^2=0\) 且 \(g\) 可逆，得 \(f^2=0\)，所以 \(u=0\)；把系数等式除以 \(g\)，同理得到 \(m/1=bg^2=1/g\)。这样构造的元素是
\[
m=\overline{-10\cdot4+3\cdot9}=\overline{11}.
\]
若另一个元素具有相同局部像，二者之差 \(h\) 满足 \(f^rh=g^sh=0\)，其中可取 \(r,s\ge1\)。理想 \((f^r,g^s)\) 为单位理想，否则包含它的极大理想将同时包含 \(f,g\)。因此 \(h=0\)，证明唯一性。
\end{solution}
